% =============================================================================
\chapter{Introdução}
\label{cap:introducao}
% =============================================================================

\section{O que é o Morphe}

O Morphe é uma plataforma para executar operações de Processamento Digital de
Sinais (PDS) em hardware e comparar o resultado com o que o computador calcula.
O aluno prepara um sinal numa janela do aplicativo, escolhe a operação e recebe
de volta duas respostas lado a lado: a que a FPGA de uma placa DE1-SoC calculou
em ponto fixo e a que o próprio computador calculou em ponto flutuante. A
diferença entre as duas é o objeto de estudo.

São cinco as operações executadas em hardware:

\begin{description}
  \item[Convolução] linear de dois sinais de até 1024 amostras cada, com saída
    de até 2047 amostras.
  \item[Filtro FIR] a mesma convolução, com a resposta ao impulso vinda de um
    projetista de filtros por especificação.
  \item[Filtro IIR] cascata de até 16 seções de segunda ordem (ordem até 32),
    também com projetista próprio.
  \item[FFT] transformada de Fourier discreta de 1024 pontos.
  \item[IFFT] a transformada inversa, no mesmo núcleo da FFT.
\end{description}

O hardware trabalha com blocos de tamanho fixo (1024 amostras), mas o
aplicativo aceita sinais de qualquer tamanho: divide o sinal em blocos, manda
cada um à placa e recompõe o resultado, com técnicas que dão a mesma resposta
que uma operação única daria (capítulo~\ref{cap:cliente}).

\section{Para que serve no ensino de PDS}

Nas disciplinas de PDS as operações são estudadas com números reais, e o
computador as calcula em ponto flutuante de 64 bits, que para quase todo fim
didático é o mesmo que calcular com números reais. Um circuito digital dedicado
não trabalha assim: usa ponto fixo, com um número pequeno de bits para a parte
fracionária, e cada multiplicação precisa ser arredondada ou truncada.

O Morphe torna essa diferença visível e mensurável. Com ele o aluno pode:

\begin{itemize}
  \item ver o erro de quantização de uma FFT em decibéis e observar que ele
    melhora 6\,dB a cada bit de faixa aproveitada (capítulo~\ref{cap:numerica});
  \item provocar uma saturação ou um estouro de acumulador e reconhecer o
    sintoma na saída;
  \item comparar um filtro IIR projetado em ponto flutuante com o mesmo filtro
    executado em ponto fixo, e entender por que o hardware usa seções de segunda
    ordem e não um polinômio de ordem alta;
  \item medir quanto tempo cada operação leva na placa e relacionar esse tempo
    com o número de operações do algoritmo.
\end{itemize}

\section{Histórico}

O Morphe nasceu como Trabalho de Conclusão de Curso de Carlos Valadão na UEFS,
\emph{Implementação de Ferramentas de Computação Numérica para o Processamento
Digital de Sinais em Hardware Reconfigurável}. A arquitetura, o núcleo de
convolução \texttt{conv1d}, a integração do núcleo de FFT, o servidor em C e o
cliente gráfico em Python são dele, e o projeto é distribuído sob a licença GNU
GPL versão~3.

A partir de setembro de 2026 a plataforma foi continuada num estágio curricular
no LabPS/UEFS, por Antonio Nicassio Santos Lima, sob supervisão do Prof.\ Armando
S.\ Sanca. O objetivo do estágio foi tornar a plataforma reproduzível a partir do
repositório e utilizável por uma turma de alunos que não conhece os detalhes de
baixo nível. A Tabela~\ref{tab:historico} resume o que mudou; o registro dia a
dia está em \arq{docs/DIARIO.md} e as correções feitas no projeto original, com
os motivos, em \arq{docs/RESSALVAS.md}.

\begin{table}[htbp]
  \centering
  \small
  \caption{O que o estágio acrescentou ao projeto original.}
  \label{tab:historico}
  \begin{tabularx}{\textwidth}{@{}lXX@{}}
    \toprule
    & \textbf{Projeto original} & \textbf{Hoje} \\
    \midrule
    Tamanho & 128 amostras na convolução e no FIR & 1024 amostras em todas as
      operações, e sinais maiores por blocos \\
    Operações & convolução, FIR, FFT & as três, mais IFFT e IIR \\
    Filtros & sem projetista por especificação & projetistas FIR e IIR por especificação \\
    Precisão & entrada da FFT sem normalização & normalização da entrada: de
      50 para 92\,dB de SNR na FFT \\
    Preparar a placa & passos manuais, conhecidos pelo autor & um comando,
      \arq{morphe-up.sh} \\
    Laboratório & uma estação, uma placa & duas placas, vários clientes ao mesmo
      tempo, instalação para a turma e Quartus em todos os PCs \\
    Reprodutibilidade & o que funcionava estava só na estação & tudo sai do
      repositório, com manifesto de integridade \\
    \bottomrule
  \end{tabularx}
\end{table}

\section{Visão geral}

A Figura~\ref{fig:visao-geral} é a figura de referência deste manual: todo o
resto detalha uma das partes dela.

\begin{figure}[htbp]
  \centering
  \begin{tikzpicture}[node distance=5mm]
  % Computadores
  \node[cliente, minimum width=30mm] (pc1) {PC do laboratório\\\scriptsize aplicativo Morphe};
  \node[cliente, minimum width=30mm, below=of pc1] (pc2) {PC do laboratório\\\scriptsize aplicativo Morphe};
  \node[cliente, minimum width=30mm, below=of pc2] (est) {Estação\\\scriptsize aplicativo + \texttt{morphe-up.sh}};

  % Rede
  \node[rede, right=14mm of pc2, minimum height=34mm, text width=18mm] (sw) {Rede do\\laboratório};

  % Placas
  \node[hps, right=16mm of sw.north east, anchor=north west, yshift=4mm, text width=24mm]
        (hps1) {HPS\\\scriptsize ARM + Linux\\\texttt{morphe\_server}};
  \node[fpga, right=10mm of hps1, text width=18mm, minimum height=16mm] (fp1) {FPGA\\\scriptsize aceleradores};
  \node[hps, right=16mm of sw.south east, anchor=south west, yshift=-4mm, text width=24mm]
        (hps2) {HPS\\\scriptsize ARM + Linux\\\texttt{morphe\_server}};
  \node[fpga, right=10mm of hps2, text width=18mm, minimum height=16mm] (fp2) {FPGA\\\scriptsize aceleradores};

  \begin{scope}[on background layer]
    \node[grupo, fit=(hps1)(fp1), label={[font=\scriptsize\bfseries]above:DE1-SoC (placa 1)}] (b1) {};
    \node[grupo, fit=(hps2)(fp2), label={[font=\scriptsize\bfseries]above:DE1-SoC (placa 2)}] (b2) {};
  \end{scope}

  \draw[<->, thick] (hps1) -- node[rotulo, above]{AXI} (fp1);
  \draw[<->, thick] (hps2) -- node[rotulo, above]{AXI} (fp2);

  % Rede
  \foreach \p in {pc1, pc2, est}
    \draw[<->, thick, cRede] (\p.east) -- (\p.east -| sw.west);
  \draw[<->, thick, cRede] (sw.east |- hps1) -- node[rotulo, above]{TCP 5000} (hps1.west);
  \draw[<->, thick, cRede] (sw.east |- hps2) -- node[rotulo, below]{TCP 5000\\SSH} (hps2.west);

  % JTAG
  \draw[->, thick, dashed, cFpga]
    (est.south) -- ++(0,-6mm) -| node[rotulo, pos=0.25, below]{USB-Blaster (JTAG): programa a FPGA e mantém o \emph{tether}} (fp2.south);
\end{tikzpicture}

  \caption{Visão geral: os computadores do laboratório falam com as placas pela
    rede (TCP, porta 5000); a estação programa a FPGA pelo USB-Blaster e sobe o
    servidor por SSH.}
  \label{fig:visao-geral}
\end{figure}

Cada \textbf{PC do laboratório} roda o aplicativo Morphe, escrito em Python. Ele
não precisa de nenhum software da Intel: fala com as placas só pela rede.

Cada \textbf{placa DE1-SoC} tem, no mesmo chip, um processador ARM rodando Linux
(o HPS, \emph{Hard Processor System}) e uma FPGA. No processador roda o
\arq{morphe_server}, que recebe as requisições pela rede na porta 5000; na FPGA
estão os aceleradores que fazem as contas. Os dois lados conversam por pontes
internas ao chip (seção~\ref{sec:pontes}).

A \textbf{estação} é o computador ligado às placas pelo cabo USB-Blaster. É dela
que se roda o \arq{morphe-up.sh}, que programa a FPGA pelo JTAG, compila e
reinicia o servidor na placa por SSH e registra a placa para os clientes. A
estação também roda o aplicativo, como qualquer outro PC.

\section{Como ler este manual}

Quem só vai \textbf{usar} a plataforma precisa do capítulo~\ref{cap:cliente}
(o aplicativo) e do capítulo~\ref{cap:laboratorio} (preparar as placas). Quem vai
\textbf{estender} a plataforma --- acrescentar uma operação, mudar um acelerador,
recompilar o hardware --- deve ler os capítulos em ordem: a
arquitetura~(\ref{cap:arquitetura}), o hardware~(\ref{cap:hardware}), o
servidor~(\ref{cap:servidor}), o protocolo~(\ref{cap:protocolo}) e a
representação numérica~(\ref{cap:numerica}) são o que se precisa saber antes de
mexer em qualquer um deles.

Convenções usadas no texto:

\begin{itemize}
  \item nomes de arquivos, comandos e sinais do hardware aparecem em
    \arq{fonte monoespaçada}; caminhos são relativos à raiz do repositório;
  \item endereços e constantes do hardware aparecem em hexadecimal, com o
    prefixo \texttt{0x};
  \item \(\mathrm{Q}m.n\) é um número em ponto fixo com sinal, \(m\) bits de
    parte inteira e \(n\) de parte fracionária (capítulo~\ref{cap:numerica});
  \item todo número medido vem com a condição em que foi medido. Um número sem
    condição é de especificação, não de medida.
\end{itemize}
